\SourceLeaf{336}
geraden Linien geht durch [den] Punkt, dessen Polarcoordinaten $\frac{1}{(12)}$, 12 sind (nemlich Distanz $\frac{1}{(12)}$, Richtung 12) und macht dort mit (12) den Winkel $02-01$ oder hat Azimuth $21-(02-01)$.

Am elegantesten wird die Construction, wenn man $(12)\times(13)$ anstatt 1 in die Zähler der Brüche setzt.

\medskip

\Needspace{13\baselineskip}\begin{center}[7.]\end{center}

Gegebene Punkte 1, 2, 3; Winkel an denselben $A, B, C$; Seiten $23=a$, $31=b$, $12=c$; $\mathfrak{A}, \mathfrak{B}, \mathfrak{C}$ vorgeschriebene Winkel an einem vierten Punkte.

$\beta$ Kreis, bestehend aus den Abschnitten

$\beta'$, wo Sehne $b$ unter Winkel $\mathfrak{B}$,

$\beta''$, wo diese Sehne unter Winkel $\mathfrak{B}\pm180^0$ erscheint.

$\gamma$ Kreis, bestehend aus den Abschnitten

$\gamma'$, wo Sehne $c$ unter Winkel $\mathfrak{C}$,

$\gamma''$, wo Sehne $c$ unter Winkel $\mathfrak{C}\pm180^0$ erscheint.

Die Kreise schneiden sich in $A$ und $D$.

In $D$ seien die durch $b$ und $c$ subtendirten Winkel $\mathfrak{B}^*, \mathfrak{C}^*$; die Sinus von $\mathfrak{B}^*-B$, $\mathfrak{C}^*-C$, $\mathfrak{B}^*+\mathfrak{C}^*-B-C$ werden alle drei positiv sein, wenn I) $D$ innerhalb des Kreises durch 1, 2, 3 liegt, alle drei hingegen negativ, wenn II) $D$ ausserhalb dieses Kreises liegt.

\Needspace{9\baselineskip}In Beziehung auf die Lage von $D$ in den Kreisen $\beta$ und $\gamma$ sind vier Fälle möglich. Es liegt $D$

1) in $\beta'$ und $\gamma'$, wo dann $\mathfrak{B}^*=\mathfrak{B}$, $\mathfrak{C}^*=\mathfrak{C}$

2) in $\beta''$ und $\gamma'$, wo $\mathfrak{B}^*=\mathfrak{B}\pm180^0$, $\mathfrak{C}^*=\mathfrak{C}$

3) in $\beta'$ und $\gamma''$, wo $\mathfrak{B}^*=\mathfrak{B}$, $\mathfrak{C}^*=\mathfrak{C}\pm180^0$

4) in $\beta''$ und $\gamma''$, wo $\mathfrak{B}^*=\mathfrak{B}\pm180^0$, $\mathfrak{C}^*=\mathfrak{C}\pm180^0$.

Im Falle 1) ist die Aufgabe physisch möglich, in den drei andern unmöglich.
